\chapter{Tổng quan về Công ty TNHH Quốc tế Khánh Sinh}

\section{Thông tin doanh nghiệp}
Công ty TNHH Quốc Tế Khánh Sinh được thành lập theo Quyết định số 3839 TC/QĐ ngày 30/12/1997 của Sở Kế hoạch và Đầu tư Thành phố Hà Nội. 
Sau khi thành lập, công ty được cấp con dấu riêng, mở tài khoản tại ngân hàng và đặt trụ sở chính tại huyện Từ Liêm, Hà Nội.\cite{company-introduction}
Công ty TNHH Quốc Tế Khánh Sinh là doanh nghiệp hoạt động trong lĩnh vực sản xuất và kinh doanh phân bón, với định hướng cung cấp các sản phẩm có chất lượng ổn định, đa dạng về chủng loại và phù hợp với nhiều loại cây trồng cũng như điều kiện thổ nhưỡng khác nhau.\cite{company-introduction} 
Trong quá trình hoạt động, công ty chú trọng đầu tư vào chất lượng sản phẩm, cải tiến kỹ thuật và nâng cao năng lực của đội ngũ cán bộ, công nhân viên. \cite{company-introduction}
Với phương châm “Chất lượng sản phẩm là sự sống còn của Công ty”, Khánh Sinh luôn hướng tới việc hoàn thiện sản phẩm, đáp ứng ngày càng tốt hơn nhu cầu của người sử dụng và từng bước củng cố vị thế của thương hiệu trên thị trường phân bón Việt Nam.\cite{company-introduction}

\begin{table}[H]
  \centering
  \caption{Thông tin khái quát về Công ty TNHH Quốc tế Khánh Sinh.\cite{company-internal-documents}}
  \label{tab:khanh-sinh-company-information}
  \begin{tabular}{p{0.31\textwidth} p{0.55\textwidth}}
    \toprule
    Nội dung & Thông tin \\
    \midrule
    Tên doanh nghiệp & Công ty TNHH Quốc tế Khánh Sinh. \\
    Tên giao dịch & Công ty Quốc tế Khánh Sinh. \\
    Tên tiếng Anh & \textit{Khanh Sinh Inter Co., Ltd.} \\
    Lĩnh vực hoạt động & Sản xuất và cung cấp các loại phân bón. \\
    Mã số thuế & 0100776325. \\
    Trụ sở chính & Miêu Nha, Tây Mỗ, Từ Liêm, Hà Nội \\
    Số điện thoại & 02438390176\\
    Website & \url{https://khanhsinh.vn}. \\
    \bottomrule
  \end{tabular}
\end{table}


\section{Lịch sử hình thành và phát triển}
\subsection{Giai đoạn trước năm 1997}
Trước năm 1997, tiền thân của Công ty TNHH Quốc Tế Khánh Sinh chỉ là một xưởng sản xuất phân bón có quy mô nhỏ. 
Diện tích nhà xưởng khi đó khoảng 400 m² với lực lượng lao động chỉ khoảng 30 người.\cite{company-internal-document} 

Trong giai đoạn này, điều kiện về cơ sở vật chất, máy móc và năng lực sản xuất còn tương đối hạn chế nên năng suất của xưởng chỉ đạt khoảng 3–4 tấn sản phẩm mỗi ngày.\cite{company-internal-document} 
Sản phẩm của cơ sở chủ yếu được cung cấp cho người nông dân tại một số huyện ngoại thành Hà Nội như Từ Liêm, Sóc Sơn và các khu vực lân cận Hà Tây. \cite{company-internal-document}
Do đó, thị trường tiêu thụ lúc bấy giờ còn mang tính địa phương và có quy mô tương đối nhỏ, phù hợp với năng lực sản xuất của cơ sở.

Mặc dù còn nhiều hạn chế về quy mô và trình độ công nghệ, giai đoạn này đã giúp doanh nghiệp tích lũy những kinh nghiệm ban đầu trong hoạt động sản xuất phân bón, xây dựng thị trường và hình thành mối quan hệ với người tiêu dùng. 
Đây cũng là nền tảng quan trọng cho quá trình thành lập và phát triển của Công ty TNHH Quốc Tế Khánh Sinh sau này.

\subsection{Giai đoạn từ năm 1997 đến năm 2014}
Năm 1997 đánh dấu bước chuyển quan trọng trong quá trình phát triển khi Công ty TNHH Quốc Tế Khánh Sinh chính thức được thành lập.\cite{company-introduction} 
Từ một cơ sở sản xuất phân bón nhỏ, công ty từng bước chuyển sang mô hình doanh nghiệp có tổ chức, đồng thời tiến hành mở rộng quy mô sản xuất và thị trường tiêu thụ.

Trong giai đoạn 1997–2014, nhu cầu sử dụng phân bón trong sản xuất nông nghiệp ngày càng tăng đã tạo điều kiện thuận lợi cho công ty phát triển. 
Để đáp ứng nhu cầu của thị trường, doanh nghiệp từng bước đầu tư thêm máy móc, thiết bị, mở rộng diện tích nhà xưởng và nâng cao năng lực sản xuất.
Đến năm 2014, doanh nghiệp đã xây dựng hai cơ sở sản xuất phân bón với tổng quy mô trên 1.200 m². \cite{company-internal-document} 
Sự gia tăng về cơ sở vật chất cho thấy hoạt động sản xuất của công ty đã có bước phát triển đáng kể so với thời kỳ đầu.

Bên cạnh việc mở rộng quy mô sản xuất, công ty cũng từng bước đa dạng hóa các dòng sản phẩm phân bón nhằm đáp ứng yêu cầu dinh dưỡng của nhiều loại cây trồng khác nhau. \cite{company-internal-document} 
Nguồn nguyên liệu phục vụ hoạt động sản xuất được mở rộng, bao gồm cả nguyên liệu trong nước và nhập khẩu. \cite{company-internal-document} 

Hoạt động phân phối sản phẩm cũng có sự thay đổi rõ rệt trong giai đoạn này. 
Từ thị trường ban đầu chủ yếu tập trung tại các huyện ngoại thành Hà Nội và khu vực lân cận, sản phẩm phân bón Khánh Sinh từng bước được đưa đến nhiều tỉnh, thành phố khác như Hà Tây, Hưng Yên, Hải Dương, Quảng Ninh và nhiều địa phương ở khu vực phía Bắc. \cite{company-internal-document} 

Công ty đồng thời chú trọng xây dựng mạng lưới nhà phân phối, đại lý và các điểm bán hàng. 
Các hoạt động khảo nghiệm sản phẩm, hội thảo kỹ thuật và hướng dẫn sử dụng phân bón cũng được triển khai nhằm giới thiệu sản phẩm cũng như hỗ trợ người nông dân lựa chọn và sử dụng phân bón phù hợp với từng loại cây trồng.\cite{company-internal-document}\cite{company-fvieu-products}

Có thể xem giai đoạn 1997–2014 là thời kỳ công ty tập trung xây dựng nền tảng và mở rộng quy mô. Trong khoảng thời gian này, Công ty TNHH Quốc Tế Khánh Sinh từng bước hoàn thiện bộ máy tổ chức, nâng cao năng lực sản xuất, phát triển danh mục sản phẩm và xây dựng hệ thống phân phối, tạo tiền đề cho giai đoạn đầu tư công nghệ và mở rộng thị trường mạnh mẽ hơn sau năm 2014.

\subsection{Giai đoạn từ năm 2014 đến nay}
Từ năm 2014, Khánh Sinh tập trung hơn vào nâng cao năng lực sản xuất, đổi mới công nghệ và đa dạng hóa sản phẩm. 
Một dấu mốc quan trọng diễn ra năm 2017 khi công ty đưa nhà máy sản xuất phân bón Fvieu vào hoạt động. 
Theo website chính thức, nhà máy sử dụng dây chuyền công nghệ Euragglo và kết hợp với nguồn nguyên liệu nhập khẩu để sản xuất các dòng phân bón Fvieu \cite{company-introduction}.

Việc đầu tư nhà máy và dây chuyền sản xuất mới thể hiện định hướng nâng cao chất lượng sản phẩm song song với mở rộng năng lực sản xuất. 
Website công ty hiện phân loại sản phẩm theo các nhóm phân bón hữu cơ, sản phẩm Fvieu, phân bón một màu, phân bón đa màu và phân đơn \cite{company-home}. 
Riêng nhóm Fvieu có nhiều công thức NPK, chẳng hạn 10-3-25, 15-15-15+TE và 13-8-12, nhằm đáp ứng yêu cầu dinh dưỡng của các nhóm cây trồng và điều kiện canh tác khác nhau \cite{company-fvieu-products}.

Nhìn chung, giai đoạn từ năm 2014 đến nay là thời kỳ doanh nghiệp đầu tư sâu hơn vào hệ thống sản xuất, công nghệ và danh mục sản phẩm. 
Trong đó, việc đưa nhà máy Fvieu với công nghệ Euragglo vào hoạt động năm 2017 là một dấu mốc tiêu biểu trong quá trình hiện đại hóa sản xuất của công ty \cite{company-introduction}.

\section{Cơ cấu tổ chức bộ máy quản trị của công ty}
Để hoạt động sản xuất kinh doanh được triển khai hiệu quả, doanh nghiệp cần xây dựng cơ cấu tổ chức quản lý phù hợp với quy mô và đặc điểm hoạt động. 
Tại Công ty TNHH Quốc tế Khánh Sinh, bộ máy quản trị được tổ chức theo mô hình chuyên chức năng, trong đó các quyết định quản lý được triển khai theo hệ thống phân cấp từ cấp lãnh đạo xuống các phòng ban và bộ phận chuyên môn \cite{company-internal-document}.

Theo mô hình này, Hội đồng quản trị thực hiện vai trò định hướng và giám sát chung đối với hoạt động của doanh nghiệp. 
Giám đốc trực tiếp điều hành hoạt động sản xuất kinh doanh, đồng thời tổ chức thực hiện các chủ trương và kế hoạch đã được phê duyệt. 
Dưới sự quản lý của Giám đốc là Phòng Hành chính, Phòng Tài chính--Kế toán, Phòng Kinh doanh, Phòng Kỹ thuật, phòng thí nghiệm kiểm định chất lượng và Bộ phận Sản xuất \cite{company-internal-document}.

Mỗi phòng ban đảm nhận chức năng, nhiệm vụ chuyên môn riêng nhưng vẫn phối hợp với các đơn vị khác để bảo đảm hoạt động chung được thực hiện thống nhất. 
Tại bộ phận sản xuất, công nhân chịu sự quản lý trực tiếp của quản đốc phân xưởng và tổ trưởng sản xuất. 
Trong khi đó, công tác hành chính, kế toán, kinh doanh và kỹ thuật do nhân viên tại các phòng ban chuyên môn tương ứng đảm nhiệm \cite{company-internal-document}.

{\renewcommand{\figurename}{Sơ đồ}
\renewcommand{\thefigure}{\arabic{figure}}
\begin{figure}[H]
  \centering
  \includegraphics[width=0.90\textwidth,height=0.74\textheight,keepaspectratio]{so_do_bo_may.png}
  \caption{Cơ cấu tổ chức bộ máy quản trị của Công ty TNHH Quốc tế Khánh Sinh.}
  \label{fig:khanh-sinh-management-model}
\end{figure}

}

Hội đồng quản trị là cấp quản lý cao nhất của công ty, có nhiệm vụ định hướng chiến lược, giám sát hoạt động sản xuất kinh doanh và quyết định các vấn đề quan trọng liên quan đến quá trình phát triển của doanh nghiệp.


Giám đốc là người trực tiếp tổ chức, quản lý và điều hành các hoạt động hàng ngày của công ty. 
Giám đốc chịu trách nhiệm triển khai các chủ trương, kế hoạch của Hội đồng quản trị và điều phối hoạt động giữa các phòng ban nhằm đảm bảo hoạt động sản xuất kinh doanh diễn ra hiệu quả.

Phòng Hành chính phụ trách các công việc liên quan đến quản lý nhân sự, văn thư, hồ sơ và các hoạt động hành chính nội bộ. 
Bộ phận này đồng thời hỗ trợ các phòng ban khác trong công tác tổ chức và đảm bảo điều kiện làm việc cho cán bộ, công nhân viên.

Phòng Tài chính - Kế toán thực hiện công tác hạch toán, quản lý thu chi, theo dõi công nợ và lập các báo cáo tài chính, kế toán của công ty. 
Ngoài ra, phòng còn tham mưu cho Ban Giám đốc về việc sử dụng và quản lý nguồn vốn phục vụ hoạt động sản xuất kinh doanh.

Phòng Kinh doanh chịu trách nhiệm tổ chức hoạt động tiêu thụ sản phẩm, chăm sóc khách hàng và duy trì quan hệ với các nhà phân phối, đại lý. 
Phòng cũng thực hiện nghiên cứu thị trường, tìm kiếm thị trường mới và xây dựng kế hoạch kinh doanh theo từng giai đoạn.

Phòng Kỹ thuật phụ trách các vấn đề chuyên môn và kỹ thuật trong quá trình sản xuất. 
Bộ phận này hướng dẫn quy trình sản xuất, theo dõi việc thực hiện các yêu cầu kỹ thuật và tham gia nghiên cứu, thử nghiệm các sản phẩm mới trên những loại cây trồng và điều kiện canh tác khác nhau.

phòng thí nghiệm kiểm định chất lượng có nhiệm vụ kiểm tra và đánh giá chất lượng nguyên liệu đầu vào, bán thành phẩm và sản phẩm hoàn thiện. 
Bộ phận này phối hợp với Phòng Kỹ thuật và Bộ phận Sản xuất nhằm kiểm soát các chỉ tiêu chất lượng, phát hiện những sai lệch trong quá trình sản xuất và góp phần đảm bảo sản phẩm trước khi đưa ra thị trường đáp ứng các yêu cầu kỹ thuật của công ty.

Bộ phận Sản xuất là lực lượng trực tiếp thực hiện quá trình sản xuất phân bón của công ty, từ tiếp nhận nguyên liệu, phối trộn đến hoàn thiện thành phẩm. 
Các phân xưởng và tổ sản xuất được quản lý bởi quản đốc và tổ trưởng nhằm đảm bảo tiến độ cũng như yêu cầu kỹ thuật của từng lô sản phẩm.
Bên cạnh hoạt động sản xuất trực tiếp, bộ phận này còn phối hợp với các đơn vị bốc xếp và vận tải trong việc lưu kho, xuất hàng và vận chuyển sản phẩm đến các địa điểm tiêu thụ. 

\section{Các nhóm sản phẩm phân bón tại Công ty TNHH Quốc tế Khánh Sinh}
Các sản phẩm phân bón của Công ty TNHH Quốc Tế Khánh Sinh được phát triển với nhiều công thức và hàm lượng dinh dưỡng khác nhau nhằm đáp ứng nhu cầu sử dụng trên nhiều nhóm cây trồng và điều kiện canh tác. 
Công ty hiện tổ chức sản xuất và kinh doanh 33 sản phẩm phân bón đã được công nhận lưu hành tại Việt Nam, căn cứ theo các quyết định công nhận phân bón lưu hành do Cục Bảo vệ thực vật, Bộ Nông nghiệp và Phát triển nông thôn ban hành năm 2023, với thời hạn lưu hành kéo dài đến năm 2028.\cite{company-internal-document}
Từ danh mục sản phẩm đang được lưu hành, có thể phân chia các sản phẩm của Khánh Sinh thành ba nhóm chính gồm phân đơn, phân NPK và phân hữu cơ - khoáng.
Mỗi nhóm có đặc điểm về thành phần dinh dưỡng và mục đích sử dụng khác nhau, qua đó giúp công ty đáp ứng tương đối đa dạng nhu cầu dinh dưỡng của cây trồng trong từng giai đoạn sinh trưởng và trên các điều kiện thổ nhưỡng khác nhau.

\subsection{Phân đơn}
Đối với nhóm phân đơn, Công ty TNHH Quốc Tế Khánh Sinh hiện có hai dòng sản phẩm chính là phân urê và phân kali clorua (KCl).

Phân urê là dòng sản phẩm cung cấp nguồn đạm cho cây trồng, trong khi kali clorua là sản phẩm cung cấp kali. 
Hai dòng phân đơn này được công ty lưu hành dưới dạng sản phẩm riêng, đồng thời cũng là những nguồn nguyên liệu quan trọng có thể được sử dụng trong quá trình phối trộn để tạo ra các công thức phân bón đa dinh dưỡng khác.

Việc duy trì các sản phẩm phân đơn giúp công ty đáp ứng nhu cầu bổ sung riêng từng nguyên tố dinh dưỡng của cây trồng, đồng thời tạo sự linh hoạt trong danh mục sản phẩm cung cấp ra thị trường.

\subsection{Phân NPK}
Phân đa dinh dưỡng là nhóm sản phẩm có số lượng và công thức tương đối đa dạng trong danh mục sản phẩm của Khánh Sinh. 
Nhóm này chủ yếu gồm các loại phân NPK và NK với nhiều tỷ lệ phối trộn khác nhau, phù hợp với nhu cầu dinh dưỡng của từng loại cây trồng và từng giai đoạn sinh trưởng.

Đối với phân NPK, công ty phát triển nhiều công thức khác nhau dựa trên tỷ lệ phối trộn giữa ba nguyên tố dinh dưỡng chính là đạm, lân và kali. 
Bên cạnh các công thức NPK cơ bản, một số sản phẩm còn được phối trộn thêm các nguyên tố trung lượng và vi lượng như lưu huỳnh (S), canxi (Ca), kẽm (Zn), bo (B) và một số nguyên tố khác. 
Việc bổ sung các thành phần này giúp công thức sản phẩm đa dạng hơn và phù hợp với yêu cầu dinh dưỡng đặc thù của từng loại cây trồng, từng loại đất và điều kiện canh tác khác nhau.

Ngoài phân NPK, công ty còn sản xuất các dòng phân NK, trong đó tỷ lệ giữa đạm và kali được điều chỉnh theo từng công thức sản phẩm. 
Sự đa dạng về tỷ lệ phối trộn cho phép công ty xây dựng nhiều dòng sản phẩm chuyên biệt, phục vụ các mục đích sử dụng khác nhau trong quá trình chăm sóc cây trồng.

Nhìn chung, nhóm phân đa dinh dưỡng thể hiện rõ định hướng đa dạng hóa sản phẩm của Khánh Sinh, không chỉ thông qua việc thay đổi tỷ lệ N-P-K mà còn thông qua việc bổ sung các nguyên tố trung và vi lượng nhằm tạo ra nhiều công thức phù hợp với nhu cầu thực tế của thị trường.

\subsection{Phân hữu cơ--khoáng}
Bên cạnh các dòng phân vô cơ, Công ty TNHH Quốc Tế Khánh Sinh còn phát triển nhóm phân hữu cơ - khoáng. 
Đây là nhóm sản phẩm kết hợp thành phần hữu cơ với các nguyên tố dinh dưỡng khoáng nhằm tạo ra sản phẩm vừa cung cấp dinh dưỡng cho cây trồng vừa bổ sung chất hữu cơ cho đất.

Các sản phẩm hữu cơ - khoáng của công ty được xây dựng với những công thức khác nhau tùy theo mục đích sử dụng và nhóm cây trồng. 
Thành phần sản phẩm có thể bao gồm chất hữu cơ kết hợp với các nguyên tố đa lượng như N, P, K và một số thành phần dinh dưỡng bổ sung khác.

Việc phát triển nhóm phân hữu cơ - khoáng giúp danh mục sản phẩm của Khánh Sinh trở nên toàn diện hơn, bên cạnh các sản phẩm phân đơn và phân đa dinh dưỡng. 
Đây cũng là nhóm sản phẩm phù hợp với định hướng nâng cao hiệu quả sử dụng phân bón đồng thời duy trì và nâng cao chất lượng đất trong quá trình canh tác lâu dài.
